\section{The autoresearch loop}

\begin{frame}{The shared greedy loop}
  \begin{center}
  \begin{tikzpicture}[node distance=0.5cm and 0.7cm]
    \node[agent]  (prop) {Propose\\\scriptsize LLM};
    \node[agent, right=of prop] (map) {Map params\\$\rightarrow$ netlist / config};
    \node[agent, right=of map] (eval) {Evaluate\\\scriptsize SPICE | LibreLane};
    \node[agent, below=of eval] (fom) {Score FoM\\\scriptsize weighted sum};
    \node[agent, left=of fom] (log) {Log to\\\cmd{results.tsv}};
    \node[agent, left=of log] (prog) {Append to\\\cmd{program.md}};
    \draw[arrow] (prop) -- (map);
    \draw[arrow] (map) -- (eval);
    \draw[arrow] (eval) -- (fom);
    \draw[arrow] (fom) -- (log);
    \draw[arrow] (log) -- (prog);
    \draw[arrow] (prog) |- (prop);
  \end{tikzpicture}
  \end{center}
  \vspace{0.3em}
  \footnotesize
  \begin{itemize}
    \item Same loop for both domains; different \term{evaluator} and
          \term{design-space} plug-ins.
    \item \cmd{program.md} -- narrative log the LLM reads back on the
          next turn. Keeps context small.
    \item \cmd{results.tsv} -- numeric history; enables plot/resume.
    \item \cmd{budget} -- hard cap on iterations. Analog: tens.
          Digital: a handful.
  \end{itemize}
  \note{%
    This is the Karpathy-autoresearch pattern applied twice. Each loop
    iteration is a proposal, a mapping step (params $\rightarrow$ a
    deck or a config), an external evaluation, a score, and a log.\par
    Why persist to markdown + TSV? Two reasons: (1) it is human-readable
    so the user can inspect the run mid-flight, (2) the LLM can re-hydrate
    context by reading only \cmd{program.md} + the last few rows of
    \cmd{results.tsv}, instead of replaying the entire prompt history
    (context windows are expensive).\par
    Design-space shape is the big difference between the two loops:
    analog is continuous in W/L and bias (we snap to grid), digital is
    discrete over knob sets (density \cmd{\{60,65,70\}}, clock \cmd{\{50,40,30\}}).
    The runner handles both via a uniform \cmd{Parameter} interface.}
\end{frame}
